\section{Forecaster Interface and Notation}
\label{app:fm-notation}

\subsection{Credit exposure}
Let $X_G(q)$ be the set of tokens issued by protocol $q$ and
$X_p(t)$ the set of tokens held by protocol $p$ at time $t$. Protocol
$p$ has a credit exposure to protocol $q$ at $t$ if
$X_p(t)\!\cap\!X_G(q)\!\neq\!\emptyset$~\citep{wu2025dexposuredatasetbenchmarksinterprotocol}.
This token-mediated dependence generalises counterparty credit risk
from traditional finance to DeFi.

\subsection{Network representation}
Over an interval $\tau\!=\!t_2\!-\!t_1$, exposures form a weighted
directed graph $G_\tau\!=\!(P_\tau,E_\tau)$. Node weights are
$w_\tau(p)\!=\!\sum_{\sigma\in X_p(t_1)\cap X_p(t_2)} v_{t_2,\sigma}$
where $v_{t_2,\sigma}$ is the USD value of token $\sigma$ at $t_2$;
edge weights aggregate non-negative token flows from $p$ to $q$
mapped to issuer $q$, as defined
in~\citet{wu2025dexposuredatasetbenchmarksinterprotocol}. We adopt
the same construction without modification.

\subsection{Forecast primitive and predicted graph}
At each decision epoch $t$ the agent queries
$\mathcal{P}_{t,h}\!\leftarrow\!\text{\fmname.Forecast}(G_t,X_t,h)$.
From $\mathcal{P}_{t,h}$ we build a representative predicted graph
via a thresholded expected-weight operator,
\begin{align}
\widehat{w}_{pq,t+h} &:= \Pr(e_{pq,t+h}\!=\!1)\cdot
                        \mathbb{E}[w_{pq,t+h}\mid e_{pq,t+h}\!=\!1],
\label{eq:expected-graph} \\
\widehat{E}_{t+h}    &:= \{(p,q) : \Pr(e_{pq,t+h}\!=\!1) > \pi_{\min}\},
\end{align}
with $\pi_{\min}\!=\!0.2$ (tuned on validation). For uncertainty
quantification we draw $M\!=\!50$ Monte Carlo samples
$\{\widetilde{G}^{(m)}_{t+h}\}$ by Bernoulli edge sampling and
Gaussian weight perturbation of $\mathcal{P}_{t,h}$.

\subsection{Hybrid prediction strategy}
In practice \fmname{} keeps all edges from $G_t$, reweights them as
$w'_{pq}\!=\!w_{pq} + \Pr(e_{pq,t+h}\!=\!1)\cdot\widehat{r}_{pq}$
where $\widehat{r}_{pq}$ is the FM residual, and additionally adds
new edges whose existence probability exceeds $\pi_{\min}$. This
preserves network structure while letting the FM adjust weights and
predict new connections.

\subsection{Architecture and training}
\label{ssec:fm-spec}
\fmname{} is built on the GraphPFN graph-tabular foundation
backbone~\citep{HollmannEtAl2025TabPFN,
eremeev2025turningtabularfoundationmodels} with a LiMiX-16M tabular
transformer encoder: $12$ layers, $4$ attention heads, hidden
dimension $64$. The same backbone serves three heads---edge
existence (binary), edge weight (regression), and node-TVL change
(regression)---trained jointly with a BCE$+$MSE composite loss. We
fine-tune one checkpoint per horizon $h\!\in\!\{1,4,8,12\}$ for
$20$ epochs with AdamW (learning rate $10^{-4}$ for heads,
$10^{-5}$ for the backbone). Fine-tuning draws on \corpusname{}
snapshots from March~2020 to January~2025: $104$ training weeks,
$12$ validation weeks (used for early stopping), and $8$ subsequent
held-out weeks for internal evaluation. All data that influences
weights or checkpoint selection precedes the frozen 2025
leaderboard split (\S\ref{app:bench-details}). Full-graph
in-context inference fits a single RTX~4090; one weekly forecast
across all four horizons completes in well under a minute.

Held-out link-prediction quality is
AUPRC$\,=\!0.978$ / $0.973$ / $0.967$ / $0.967$ and
AUROC$\,=\!0.996$ / $0.995$ / $0.993$ / $0.993$ at $h\!=\!1$ / $4$
/ $8$ / $12$ respectively; a frozen GraphPFN baseline (no
fine-tuning) achieves AUPRC $0.936$--$0.940$ at the same horizons,
isolating the contribution of domain fine-tuning. These intrinsic
forecaster numbers do not appear in the main text because the
paper's contribution is the operating point of the full pipeline,
not the FM's headline accuracy.
